\section{Results}\label{sec:results}

\paragraph{Decision rules, fixed before the results.} A measure is usable if it scores most videos, a topic's own
video scores higher than another topic's video, and two measurements of the same videos agree. The learner repair
helps if it beats the generic repair of the same size; finding the weak section matters if it also beats rewriting
every section. A pipeline feature helps if it raises these differences without hurting the layout. All numbers are
means over topics on a 0--100 scale; paired bootstrap intervals and Wilcoxon tests are kept with the result files.
We call a difference \emph{beyond noise} only if it passes two tests. First, it must survive randomly swapping the two
versions' answers question by question. Second, its sign must hold across topics. The student is not fully
repeatable. We re-asked \NMRepeatN{} questions on identical frames, and its confidence in the right answer moved on
\NMRepeatMoved{} of them, a few times by up to half the scale. On one concept, repeat noise alone moves M9 by up to
\NMFloorOne{} points (95\%); on the mean of all concepts, by up to \NMFloorAll{}.

\paragraph{The ladder on MMMC} (Table~\ref{tab:ladder_mmmc}). Our analysis yields four insights.
(i) \textbf{The student already knows MMMC.} Forced to choose, it answers every question right with no video on
\LadMMMCBlindAllRight{} of \LadMMMCTopics{} pilot topics. There is nothing left to learn, so every gain measure
(M3--M9) scores \LadMMMCDefinedNVG{} videos.
(ii) \textbf{TeachQuiz sees the topic, not the repair.} With answers matched in order (M1), a topic's own video beats
another topic's video by \LadMMMCSepOne{} points, but the learner repair beats the base video by only
\LadMMMCLBOne{}, and \LadMMMCTiesOne{} of \LadMMMCTopics{} topics are ties. With forced choice (M2), all
\LadMMMCTiesTwo{} topics are ties.
(iii) \textbf{The raw scores sit at the prior ceiling.} Without normalisation, confidence in the right answer is
\LadMMMCSBlind{} with no video and \LadMMMCSVisual{} with the dense video. Another topic's video lowers it further,
which is all that the own-versus-other difference of R9 (\LadMMMCSepRawNine{}) measures.
(iv) \textbf{A smaller student does not change the picture.} With \texttt{qwen2.5vl:3b}, confidence with no video is
\LadSmallSBlind{}, so the gain on 3 frames per section (M9) scores \LadSmallDefinedNine{} videos
(Table~\ref{tab:ladder_small}). Its own video beats another topic's by \LadSmallSepNine{} points, but the learner
repair scores \LadSmallLBNine{} against the base video: each topic's gain rests on one or two questions. With so
little room left, a learner-versus-base difference would have to reach about \NMSmallMDE{} points before this design
could detect it.
On MMMC, with this student, the repair can neither succeed nor fail; the pilot's null result is not evidence
against it.

\begin{table*}[t]\centering\footnotesize
\caption{A smaller student (qwen2.5vl:3b) on the MMMC pilot videos, 3 frames per section only; measures it could not score are left out. 12 topics, 48 videos. Each row adds one change to the row above; points on a 0--100 scale. Scored: videos that get a score (-- : the student already knew everything). Own $-$ other: a topic's video minus another topic's video on the same quiz. Agreement: rank correlation of two measurements of the same videos. Learner $-$ \ldots: mean change over topics. Ties: topics where learner and base score the same.}\label{tab:ladder_small}
\setlength{\tabcolsep}{4pt}\begin{tabular}{@{}llrrrrrrr@{}}\toprule & & \multicolumn{3}{c}{Does it measure teaching?} & \multicolumn{4}{c}{Does the learner repair help?}\\
\cmidrule(lr){3-5}\cmidrule(lr){6-9}
& Measure & Scored & \hd{Own $-$\\other ($\uparrow$)} & \hd{Agree-\\ment ($\uparrow$)} & \hd{Learner\\$-$ base} & \hd{Learner\\$-$ generic} & \hd{Learner $-$\\every section} & Ties ($\downarrow$)\\\midrule
M0 & TeachQuiz (Code2Video) & 48/48 & -- & -- & +1.7 & $-$2.2 & $-$4.5 & 7/12 \\
M2 & + must choose (no NULL) & 48/48 & +10.0 & 0.65 & $-$3.3 & +1.7 & +3.3 & 9/12 \\
M3 & + gain over what it knew & 24/48 & +16.7 & 0.72 & +25.0 & +8.3 & +16.7 & 4/12 \\
M9 & M7 on 3 frames per section & 36/48 & +37.9 & -- & $-$23.7 & $-$0.6 & +5.9 & 1/12 \\
R9 & 3 frames minus no video & 48/48 & +6.4 & 0.70 & $-$2.9 & +0.2 & +1.2 & 1/12 \\
\bottomrule\end{tabular}
\end{table*}


\paragraph{The ladder on NOVEL-12} (Table~\ref{tab:ladder_novel}). On the \LadNovelTopics{} admitted concepts the
student does not know the answers in advance (confidence \LadNovelSBlind{} with no video), so the gain measures give
scores. Three insights follow.
(i) \textbf{The videos teach invented concepts.} The gain on 3 frames per section (M9) scores
\LadNovelDefinedNine{} base videos. It averages \NMGain{} points and is positive on \NMGainPos{} of \NMGainN{}
concepts, which is beyond noise.
(ii) \textbf{What the student learns is the topic of the video.} A topic's own video beats another topic's video by
\LadNovelSepNine{} points of M9, on \NMSpecPos{} of \NMSpecN{} concepts. This is also beyond noise.
(iii) \textbf{Three sharp frames per section teach more than the dense video, on these videos.} Confidence on the
right answer is \LadNovelSFrames{} with 3 frames per section and \LadNovelSVisual{} with the dense video. D-VGAIN on
the dense video (M7) scores \LadNovelDefinedSeven{} videos with \LadNovelSepSeven{} points of separation. M9 is
\NMFramesDense{} points above the dense video's gain, on \NMFramesDensePos{} of \NMFramesDenseN{} concepts. The
difference is beyond noise and holds across concepts, but only just: its interval barely excludes zero. Two things
change together here: the 3 frames are larger (768 pixels wide, against 512) and far fewer. So we cannot yet say
which of the two matters.

A single concept's score is fragile. It rests on 8 questions, and a concept the student half knows leaves little room
to learn. On \emph{Sorrel counts} (confidence 66.8 with no video), one changed answer in one option order moves M9 by
about 19 points. The only negative gain in (i) is this concept's, and it is within noise.

\paragraph{Video quality} (Tables~\ref{tab:quality} and~\ref{tab:layout_study}). \AdmIncompleteBefore{} of
\AdmTotal{} NOVEL-12 base videos had silently dropped a section that failed to render. After one fresh attempt,
\AdmComplete{} were complete and enter the study. Claude (\texttt{claude-opus-5-5}), seeing only code, frames and the
automatic reports, fixed \ClaudeFaults{} visible faults in their \ClaudeSections{} sections (\ClaudeOverlap{} overlaps,
\ClaudeOther{} objects missing or mistimed, \ClaudeOffScreen{} off screen, \ClaudeLeftover{} leftovers). By every
check, the videos got cleaner. Lasting conflicts fell from \LSOursConflicts{} to \LSClaudeConflicts{}, and broken
attachments fell from \LSOursAttach{} to \LSClaudeAttach{}. Yet MVQS rose only from \LSOursMvqs{} to
\LSClaudeMvqs{}. It can even fall when a fix is right. On one \emph{Glint numbers} section, a shift had left two
highlight boxes over empty space, and Claude moved them back onto the digits they mark. MVQS fell from 88.4 to 67.0,
because the boxes now touch the digits. Our layout repair had inserted \ShiftsIn{} \texttt{shift} calls; Claude removed \ShiftsRemoved{}, and
\ShiftNotes{} of its notes name a shift as the cause of a fault. By the rule written before scoring, layout changes
learning if Claude-fixed beats ours by at least 5 points of M9 on average and on at least 7 of 10 concepts. The
difference is \LSDiff{} points with \LSWins{} of \LSN{}, so the rule is \LSGate{}. On noise alone, the rule
passes in \NMRuleFPR\% of random swaps. This is a weak null result. With 10 concepts of 8 questions each, the study
can only detect a layout effect of about \NMLayoutMDE{} points or more, far above the rule's bar of 5. No large
layout effect appeared; smaller effects remain possible. Single concepts moved both ways, from
$-25.9$ (\emph{Mirthful triangles}) to $+62.7$ (\emph{Sorrel counts}). Once we allow for testing ten concepts, none of
these moves is beyond noise. Sorrel's gain on 3 frames also reverses on the dense video ($-19.7$).

\paragraph{Pipeline ablation.} The stages F0--F5
(\S\ref{sec:ablation}) are built and their decision rules fixed; their runs, about 14 hours on the shared machine,
continue after this draft.
